\begin{ExerciseBlock}
\ExerciseHeading{04}{Youngscher Doppelspalt}{geometrical_and_wave_optics_double_slit}
Eine ebene elektromagnetische Welle der Wellenlänge $\lambda=\qty{632.8}{\nano\metre}$ trifft senkrecht auf einen Doppelspalt mit Spaltmittenabstand $d=\qty{10}{\micro\metre}$. Die Spalte sind so schmal, dass ihre Einzelspalthülle vernachlässigt werden kann. In der Fernfeldnäherung verlaufen die beiden Strahlen zum selben Schirmpunkt nahezu parallel.
\begin{center}
\CachedTikzFigure{content/resources/chapters/geometrical_and_wave_optics/tikz/young_double_slit.tex}
\end{center}
\begin{enumerate}
\item Berechnen Sie den Gangunterschied $\Delta s$ in Abhängigkeit vom Abstrahlwinkel $\alpha$.
\item Für welche Winkel $\alpha_m$ treten Intensitätsmaxima auf? Berechnen Sie für $m=1,2,3$ die Winkel und die Winkelabstände benachbarter Maxima. Was fällt Ihnen auf?
\item Welche höchste Ordnung tritt hinter dem Doppelspalt auf?
\item Berechnen Sie die Winkel der beiden höchsten positiven Ordnungen und ihren Abstand. Vergleichen Sie mit Teilaufgabe b).
\end{enumerate}
\end{ExerciseBlock}
\begin{Solution}
\begin{RemarkBox}[title=Korrektur zur handschriftlichen Lösung]
Für die vorgegebene Wellenlänge $\lambda=\qty{632.8}{\nano\metre}$ gilt
\[
 m_\mathrm{max}=\left\lfloor\frac{d}{\lambda}\right\rfloor=15.
\]
Die höchsten beiden positiven Ordnungen liegen bei $\alpha_{14}\approx\ang{62.36}$ und $\alpha_{15}\approx\ang{71.66}$. Ihr Abstand beträgt rund \ang{9.29}.
\end{RemarkBox}
\end{Solution}
\UseExerciseSolution{geometrical_and_wave_optics_double_slit}
